
\section{教育背景}
\datedsubsection{\textbf{麻省理工学院（MIT）}, 美国}{2018 -- 2022}
\begin{itemize}
\item \textbf{土木与环境工程系，交通科学博士}
\item GPA: \textbf{5.0/5.0}，15门主课涵盖计算机、运筹学、概率统计、经济学等多个方向，其中6门课程获得A+，超额完成项目课程要求（累计229学分，项目毕业要求120学分），仅用四年博士毕业
\item 获得2021年 MIT UPS 博士奖学金（约9万美金，每年奖励给一名交通、物流、优化领域最优秀的博士生），连续2年获得院系提名申请MIT William Asbjornsen 奖学金（每年奖励给一名MIT科学与工程领域最优秀的博士生，每个院系每年提名一人）
\item 导师: Jinhua Zhao, Haris N. Koutsopoulos
\item 答辩委员会成员: Jinhua Zhao, Haris N. Koutsopoulos, Max Zuo-Jun Shen, Cathy Wu
\item 博士论文获: COTA最佳博士论文奖，MIT Dan \& Eva Roos Thesis Prize, Global Urban Rail Transit Best Dissertation Award
\end{itemize}

\datedsubsection{\textbf{麻省理工学院（MIT）}, 美国}{2018 -- 2020}
\begin{itemize}
\item \textbf{土木与环境工程系，交通科学硕士}
\item \textbf{电子工程与计算机科学系，计算机科学硕士}
\item GPA: \textbf{5.0/5.0}，两年完成双硕士学位（通常需三年）
\item 导师: Jinhua Zhao, Haris N. Koutsopoulos, Patrick Jaillet
\end{itemize}

\datedsubsection{\textbf{清华大学}, 北京}{2014 -- 2018}
\begin{itemize}
\item \textbf{土木工程系，土木工程学士}
\item \textbf{经管学院，管理学学士（双学位）}
\item GPA: \textbf{93/100}， 排名 \textbf{1/105}. 获得2017年清华大学本科生特等奖学金，清华大学优秀毕业生，优秀毕业论文；2018年土木水利学院毕业典礼本科生发言代表。
\item 导师: 李瑞敏
\end{itemize}
